\chapter{Conclusiones y trabajo futuro}
\label{cap:conclusiones}

\section{Respuesta a las preguntas de investigación}
\textbf{RQ1.} El prototipo ofrece al votante un recibo y publica evidencias de
presencia e integridad que pueden explicarse sin revelar la opción escogida.
La comprensibilidad efectiva no queda demostrada hasta realizar la evaluación
con usuarios; tampoco se ofrece verificabilidad extremo a extremo.

\textbf{RQ2.} La separación entre voter-access y ballot-box, la ausencia de
identificadores y marcas temporales por voto y el token de un solo uso
constituyen los mecanismos implementados para controlar unicidad y reducir la
vinculación directa. La protección tiene límites ante un operador que controle
ambas bases o los dispositivos.

\textbf{RQ3.} ISO/IEC 25010:2023 proporciona el marco para organizar los RNF y
la evaluación. La suite automatizada aporta evidencia inicial de reglas
seleccionadas y cuatro integraciones de migraciones y esquema en PostgreSQL.
La integración de repositorios y flujos, la accesibilidad, la carga y la
usabilidad siguen pendientes, lo que impide una valoración completa de calidad.

\section{Cumplimiento de objetivos}
OBJ-1 y OE1 cuentan con arquitectura multi-tenant, organizaciones, roles y
aislamiento por organización implementados. OBJ-2 y OE2 cuentan con censos,
procesos, ciclo de vida y API documentada. OBJ-3 y OE3 cuentan con flujo de
participación, secreto por separación de datos y recibo; la accesibilidad y
comprensión requieren validación pendiente. OBJ-4 y OE4 cuentan con auditoría,
escrutinio, huellas y resultados; su evaluación de calidad sigue incompleta.
Por tanto, el cumplimiento funcional implementado no equivale a aceptación
experimental de todos los objetivos. Las limitaciones y el trabajo futuro de
la sección~\ref{sec:limitaciones-trabajo-futuro} recogen los protocolos
pendientes.

\section{Limitaciones y trabajo futuro}
\label{sec:limitaciones-trabajo-futuro}
Quedan por ejecutar la integración de repositorios y flujos con PostgreSQL,
las pruebas E2E del stack completo, la revisión manual de accesibilidad con
teclado y lector de pantalla (H1a: parte automática favorable; parte manual no realizada; la parte
automática se midió en el flujo del votante tras corregir las incidencias de
la papeleta y del recibo), tres ejecuciones medidas de
carga con k6 (H1c) y la
sesión con cinco participantes votantes (H1b), que es la evaluación
prioritaria; el consentimiento, el guion y la hoja de registro se conservan en
\texttt{docs/evaluation/usabilidad/}. La sesión de RNF-4 con cinco
administradores se realizó el 03/11/2026: los cinco completaron la tarea sin
ayuda de otras personas, pero solo dos no tenían formación previa y los
tiempos se estimaron después de la sesión, por lo que RNF-4 queda como no
demostrado y su resultado solo es indicativo. Cerrarlo requiere repetir la
sesión con administradores sin formación previa y tiempos cronometrados. Si alguna de estas
evaluaciones no se completa antes de la entrega, su resultado figura como
PENDIENTE en la tabla~\ref{tab:estado-hipotesis} y constituye una limitación
de este trabajo. También se debe ampliar la cobertura de código.
Las propuestas registradas en \texttt{docs/ai/IMPROVEMENTS.md} orientan
mejoras posteriores. P9 contempla RF-11, RF-22 y RF-23 fuera del MVP actual.
El modelo de sustitución del voto de la especificación v1.0 (estado
«Sustituido» del voto y del recibo, RI-6) se conserva descrito como referencia
de diseño para RF-11, pero no es operativo en esta versión: la urna registra un
único voto aceptado por token de emisión y nunca genera ese estado. Activarlo
requiere implementar la lógica de sustitución y validarla con pruebas de
concurrencia y unicidad.
D-03 y D-05 fueron confirmadas por el tutor el 30/10/2026 mediante correo
electrónico.
